\section{SENL (Sistemas de Ecuaciones no Lineales)}
\subsection{Newton-Raphson}
El método consigue una de las raíces del sistema de ecuaciones, habría que iterarlo varias veces para obtener todas.
El Jacobiano tiene que ser evaluado en un punto.
Fórmula iterativa desarrollada:
\begin{equation}
\begin{pmatrix}
    x_{k+1}\\y_{k+1}
\end{pmatrix}=\begin{pmatrix}
    x_k\\y_k
\end{pmatrix}-\begin{pmatrix}
  \frac{\partial f_1(x_k,y_k)}{\partial x_k} & \frac{\partial f_1(x_k,y_k)}{\partial y_k} \\ 
  \frac{\partial f_2(x_k,y_k)}{\partial x_k} & \frac{\partial f_2(x_k,y_k)}{\partial y_k} 
\end{pmatrix}^{-1}\times \begin{pmatrix}
    f_1(x_k,y_k)\\f_2(x_k,y_k)
\end{pmatrix}
\end{equation}
\begin{equation}
\vec{x}_{k+1}=\vec{x}_k-\vec{J}(\vec{x}_k)^{-1}\times\vec{F}(\vec{x}_k)
\end{equation}

\subsection{Método Punto Fijo}
El método de punto fijo es igual que el método de Newton Raphson, la diferencia está en el jacobiano. En el método de punto fijo se elije una matriz $A$ a conveniencia.
Se puede tomar la matriz jacobiana y mantener la matriz evaluada en la semilla para todas las iteraciones, para no tener que inventar una y chequear convergencia.
\begin{equation}
\vec{x}_{k+1}=\vec{x}_k-\vec{J}(\vec{x}_k)^{-1}\times\vec{F}(\vec{x}_k)
\end{equation}

\subsection{Cálculo de inversa de una matriz}
$A^{-1}=\frac1{\det(A)}\times \text{Adj}(A)$

$\text{Adj}(A)=\text{Cof}(A)^T$

$C_{ij}=(-1)^{i+j} \times \det(\text{de la matriz que queda al eliminar la fila $i$ y la col $j$})$
La matriz de cofactores es del mismo tamaño que $A$ y cada elemento es el cofactor $ij$.
Si es un solo elemento el que queda, va ese elem directamente como determinante de la matriz sobrante.